\color{uniblue}\subsection[Выбор уставок защиты от перегрузки]{Выбор уставок защиты от перегрузки}
\color{black}

\begin{enumerate}
\item
По выражению (\ref{eq:zp1}) определяется первичный ток срабатывания защиты от перегрузки $I_{\textsl{с.з.ЗП}}$:
\begin{equation*}
    I_{\textsl{с.з.ЗП}} = \frac{K_{\textsl{отс}}}{K_{\textsl{в}}} \cdot I_{\textsl{ном.ст.}} = \frac{1,05}{0,95} \cdot 250 = 276\; (\textsl{А}),
\end{equation*}

Уставка задается в относительных единицах от номинального тока аналогового входа устройства, для этого полученное значение пересчитывается по следующему выражению:
\begin{equation*}
    I_{\textsl{ср}} = \frac{I_{\textsl{с.з.ЗП}}}{I_{\textsl{ном.вх.уст}}} \cdot \frac{I_{\textsl{ном.вт.ТТ}}}{I_{\textsl{ном.пер.ТТ}}} = \frac{276}{5} \cdot \frac{5}{400} = 0,7\; (\textsl{о.е.}),
\end{equation*}
\begin{eqexpl}[15mm]
    \item{$I_{\textsl{ном.вх.уст}}$} номинальный ток аналоговых входов устройства, 1 или 5 А;
    \item{$I_{\textsl{ном.вт.ТТ}}$} номинальный вторичный ток ТТ, А;
    \item{$I_{\textsl{ном.пер.ТТ}}$} номинальный первичный ток ТТ, А.
\end{eqexpl}

\item
Параметр \textbf{$I_{\textsl{ср}}$} <<Ток срабатывания>> функции ЗП принимается \textbf{равным 0,7}.
\item
Параметр \textbf{$T_{\textsl{ср}}$} <<Выдержка времени срабатывания>> функции ЗП принимается \textbf{равным 10 c}.

\end{enumerate}